\documentclass[12pt,a4paper]{article}
\usepackage[utf8]{inputenc}
\usepackage[margin=1in]{geometry}
\usepackage{graphicx}
\usepackage{listings}
\usepackage{xcolor}
\usepackage{hyperref}
\usepackage{booktabs}
\usepackage{longtable}
\usepackage{fancyhdr}
\usepackage{titlesec}

% Code listing style
\lstset{
    basicstyle=\ttfamily\small,
    breaklines=true,
    frame=single,
    backgroundcolor=\color{gray!10},
    keywordstyle=\color{blue},
    commentstyle=\color{green!60!black},
    stringstyle=\color{red},
    numbers=none,
    showstringspaces=false
}

\pagestyle{fancy}
\fancyhf{}
\rhead{CMPE 300 - Project 2}
\lhead{MPI-Based Parallel NLP System}
\rfoot{Page \thepage}

\title{
    \textbf{CMPE 300: Analysis of Algorithms} \\
    \vspace{0.5cm}
    \Large Project 2: MPI-Based Parallel NLP System \\
    \vspace{1cm}
}

\author{
    \textbf{Student 1 Name} \\
    Student ID: XXXXXXXXX \\
    \and
    \textbf{Student 2 Name} \\
    Student ID: XXXXXXXXX
}

\date{December 2025}

\begin{document}

\maketitle
\thispagestyle{empty}
\newpage

\tableofcontents
\newpage

%==============================================================================
\section{Program Description}
%==============================================================================

In this project, we implemented an MPI-based parallel Natural Language Processing (NLP) system. The goal was to perform text preprocessing and frequency analysis on input text files. The system applies lowercasing, punctuation removal, and stopword removal, then computes term-frequency (TF) and document-frequency (DF) for a given vocabulary.

We used only \texttt{MPI\_Send} and \texttt{MPI\_Recv} for communication, avoiding collective operations as specified.

\subsection{NLP Operations}

We implemented the following operations:
\begin{itemize}
    \item \textbf{Lowercasing}: Using Python's \texttt{str.lower()}.
    \item \textbf{Punctuation Removal}: Using \texttt{str.translate()} with \texttt{string.punctuation}.
    \item \textbf{Stopword Removal}: Filtering out words present in the stopword list.
    \item \textbf{TF Counting}: Counting total word occurrences.
    \item \textbf{DF Counting}: Counting the number of sentences containing each word.
\end{itemize}

\subsection{Challenges}

One of the main challenges was handling data types correctly, as some functions receive raw strings while others receive lists of words. We solved this by checking types at runtime.

Another significant challenge was preventing deadlocks in Pattern 4 during the data exchange phase. We implemented an asymmetric communication strategy where odd-ranked processes send first and even-ranked processes receive first.

We also had to ensure proper synchronization in the pipeline patterns (2 and 3) by propagating a termination signal (\texttt{None}) through the stages. Finally, distributing data evenly among workers required careful calculation of chunk sizes to handle edge cases.

%==============================================================================
\section{Pattern Implementations}
%==============================================================================

\subsection{Pattern 1: End-to-End Processing}

This pattern requires at least 2 processes. The manager (rank 0) loads the data and splits the text into $n-1$ equal chunks. Each worker receives a chunk, processes it fully (lowercasing, punctuation removal, stopword filtering), and computes the TF counts. Finally, the workers send their local counts back to the manager, which sums them up.

\subsection{Pattern 2: Linear Pipeline}

This pattern uses exactly 5 processes. We set up a pipeline where data flows from one worker to the next.
\begin{itemize}
    \item Rank 1: Lowercases the text.
    \item Rank 2: Removes punctuation.
    \item Rank 3: Removes stopwords.
    \item Rank 4: Counts term frequencies.
\end{itemize}
The manager sends the text in small chunks to keep the pipeline busy. We used a \texttt{None} signal to tell the workers when to stop.

\subsection{Pattern 3: Parallel Pipelines}

This pattern runs multiple pipelines at once. It requires $1 + 4k$ processes. We split the text into large chunks for each pipeline, and then split those into smaller chunks for the workers inside the pipeline. This combines the benefits of parallel processing and pipelining.

\subsection{Pattern 4: Task Parallelism}

This pattern requires $1 + 2k$ processes. First, the manager distributes the text to all workers for preprocessing. Then, workers pair up (1 with 2, 3 with 4, etc.) and exchange their data so they both have the same set.
After exchanging, they split the work: odd ranks calculate TF, and even ranks calculate DF. To avoid deadlocks during the exchange, we made odd ranks send first and even ranks receive first.

%==============================================================================
\section{Test Cases}
%==============================================================================

\subsection{Test Case Specifications}

\begin{table}[h]
\centering
\begin{tabular}{|c|c|c|c|l|}
\hline
\textbf{Test} & \textbf{Sentences} & \textbf{Vocab Size} & \textbf{Stopwords} & \textbf{Theme} \\
\hline
1 & 20 & 5 & 4 & English Proverbs \\
2 & 40 & 8 & 6 & Technology/Computing \\
3 & 60 & 10 & 7 & Nature/Wildlife \\
4 & 40 & 11 & 9 & Algorithms/CS \\
5 & 100 & 15 & 9 & Turkey/Culture \\
\hline
\end{tabular}
\caption{Test Case Specifications}
\end{table}

%------------------------------------------------------------------------------
\subsection{Test Case 1 Results}
%------------------------------------------------------------------------------

\subsubsection{Pattern 1 (4 processes)}
\begin{lstlisting}
Command: mpiexec -n 4 python solution.py --text test_cases/text_1.txt --vocab test_cases/vocab_1.txt --stopwords test_cases/stopwords_1.txt --pattern 1

Output:
Running Pattern 1 with 4 processes.
Term-Frequency (TF) Result:
bird: 1
dog: 1
fox: 1
power: 2
time: 1
\end{lstlisting}

\subsubsection{Pattern 2 (5 processes)}
\begin{lstlisting}
Command: mpiexec -n 5 python solution.py --text test_cases/text_1.txt --vocab test_cases/vocab_1.txt --stopwords test_cases/stopwords_1.txt --pattern 2

Output:
Running Pattern 2 with 5 processes.
Term-Frequency (TF) Result:
bird: 1
dog: 1
fox: 1
power: 2
time: 1
\end{lstlisting}

\subsubsection{Pattern 3 (9 processes)}
\begin{lstlisting}
Command: mpiexec -n 9 python solution.py --text test_cases/text_1.txt --vocab test_cases/vocab_1.txt --stopwords test_cases/stopwords_1.txt --pattern 3

Output:
Running Pattern 3 with 9 processes.
Term-Frequency (TF) Result:
bird: 1
dog: 1
fox: 1
power: 2
time: 1
\end{lstlisting}

\subsubsection{Pattern 4 (5 processes)}
\begin{lstlisting}
Command: mpiexec -n 5 python solution.py --text test_cases/text_1.txt --vocab test_cases/vocab_1.txt --stopwords test_cases/stopwords_1.txt --pattern 4

Output:
Running Pattern 4 with 5 processes.
Term-Frequency (TF) Result:
bird: 1
dog: 1
fox: 1
power: 2
time: 1

Document-Frequency (DF) Result:
bird: 1
dog: 1
fox: 1
power: 1
time: 1
\end{lstlisting}

%------------------------------------------------------------------------------
\subsection{Test Case 2 Results}
%------------------------------------------------------------------------------

\subsubsection{Pattern 1 (4 processes)}
\begin{lstlisting}
Command: mpiexec -n 4 python solution.py --text test_cases/text_2.txt --vocab test_cases/vocab_2.txt --stopwords test_cases/stopwords_2.txt --pattern 1

Output:
Running Pattern 1 with 4 processes.
Term-Frequency (TF) Result:
code: 2
computing: 4
data: 6
learning: 2
performance: 3
security: 1
software: 5
systems: 4
\end{lstlisting}

\subsubsection{Pattern 2 (5 processes)}
\begin{lstlisting}
Command: mpiexec -n 5 python solution.py --text test_cases/text_2.txt --vocab test_cases/vocab_2.txt --stopwords test_cases/stopwords_2.txt --pattern 2

Output:
Running Pattern 2 with 5 processes.
Term-Frequency (TF) Result:
code: 2
computing: 4
data: 6
learning: 2
performance: 3
security: 1
software: 5
systems: 4
\end{lstlisting}

\subsubsection{Pattern 3 (9 processes)}
\begin{lstlisting}
Command: mpiexec -n 9 python solution.py --text test_cases/text_2.txt --vocab test_cases/vocab_2.txt --stopwords test_cases/stopwords_2.txt --pattern 3

Output:
Running Pattern 3 with 9 processes.
Term-Frequency (TF) Result:
code: 2
computing: 4
data: 6
learning: 2
performance: 3
security: 1
software: 5
systems: 4
\end{lstlisting}

\subsubsection{Pattern 4 (5 processes)}
\begin{lstlisting}
Command: mpiexec -n 5 python solution.py --text test_cases/text_2.txt --vocab test_cases/vocab_2.txt --stopwords test_cases/stopwords_2.txt --pattern 4

Output:
Running Pattern 4 with 5 processes.
Term-Frequency (TF) Result:
code: 2
computing: 4
data: 6
learning: 2
performance: 3
security: 1
software: 5
systems: 4

Document-Frequency (DF) Result:
code: 2
computing: 4
data: 6
learning: 2
performance: 3
security: 1
software: 5
systems: 4
\end{lstlisting}

%------------------------------------------------------------------------------
\subsection{Test Case 3 Results}
%------------------------------------------------------------------------------

\subsubsection{Pattern 1 (4 processes)}
\begin{lstlisting}
Command: mpiexec -n 4 python solution.py --text test_cases/text_3.txt --vocab test_cases/vocab_3.txt --stopwords test_cases/stopwords_3.txt --pattern 1

Output:
Running Pattern 1 with 4 processes.
Term-Frequency (TF) Result:
birds: 2
forest: 1
night: 2
ocean: 2
sky: 2
summer: 1
sun: 1
trees: 4
water: 2
winter: 2
\end{lstlisting}

\subsubsection{Pattern 2 (5 processes)}
\begin{lstlisting}
Command: mpiexec -n 5 python solution.py --text test_cases/text_3.txt --vocab test_cases/vocab_3.txt --stopwords test_cases/stopwords_3.txt --pattern 2

Output:
Running Pattern 2 with 5 processes.
Term-Frequency (TF) Result:
birds: 2
forest: 1
night: 2
ocean: 2
sky: 2
summer: 1
sun: 1
trees: 4
water: 2
winter: 2
\end{lstlisting}

\subsubsection{Pattern 3 (9 processes)}
\begin{lstlisting}
Command: mpiexec -n 9 python solution.py --text test_cases/text_3.txt --vocab test_cases/vocab_3.txt --stopwords test_cases/stopwords_3.txt --pattern 3

Output:
Running Pattern 3 with 9 processes.
Term-Frequency (TF) Result:
birds: 2
forest: 1
night: 2
ocean: 2
sky: 2
summer: 1
sun: 1
trees: 4
water: 2
winter: 2
\end{lstlisting}

\subsubsection{Pattern 4 (5 processes)}
\begin{lstlisting}
Command: mpiexec -n 5 python solution.py --text test_cases/text_3.txt --vocab test_cases/vocab_3.txt --stopwords test_cases/stopwords_3.txt --pattern 4

Output:
Running Pattern 4 with 5 processes.
Term-Frequency (TF) Result:
birds: 2
forest: 1
night: 2
ocean: 2
sky: 2
summer: 1
sun: 1
trees: 4
water: 2
winter: 2

Document-Frequency (DF) Result:
birds: 2
forest: 1
night: 2
ocean: 2
sky: 2
summer: 1
sun: 1
trees: 4
water: 2
winter: 2
\end{lstlisting}

%------------------------------------------------------------------------------
\subsection{Test Case 4 Results}
%------------------------------------------------------------------------------

\subsubsection{Pattern 1 (4 processes)}
\begin{lstlisting}
Command: mpiexec -n 4 python solution.py --text test_cases/text_4.txt --vocab test_cases/vocab_4.txt --stopwords test_cases/stopwords_4.txt --pattern 1

Output:
Running Pattern 1 with 4 processes.
Term-Frequency (TF) Result:
algorithms: 33
complexity: 3
data: 7
efficient: 1
memory: 2
patterns: 2
problems: 4
search: 2
solutions: 2
sort: 10
time: 3
\end{lstlisting}

\subsubsection{Pattern 2 (5 processes)}
\begin{lstlisting}
Command: mpiexec -n 5 python solution.py --text test_cases/text_4.txt --vocab test_cases/vocab_4.txt --stopwords test_cases/stopwords_4.txt --pattern 2

Output:
Running Pattern 2 with 5 processes.
Term-Frequency (TF) Result:
algorithms: 33
complexity: 3
data: 7
efficient: 1
memory: 2
patterns: 2
problems: 4
search: 2
solutions: 2
sort: 10
time: 3
\end{lstlisting}

\subsubsection{Pattern 3 (9 processes)}
\begin{lstlisting}
Command: mpiexec -n 9 python solution.py --text test_cases/text_4.txt --vocab test_cases/vocab_4.txt --stopwords test_cases/stopwords_4.txt --pattern 3

Output:
Running Pattern 3 with 9 processes.
Term-Frequency (TF) Result:
algorithms: 33
complexity: 3
data: 7
efficient: 1
memory: 2
patterns: 2
problems: 4
search: 2
solutions: 2
sort: 10
time: 3
\end{lstlisting}

\subsubsection{Pattern 4 (5 processes)}
\begin{lstlisting}
Command: mpiexec -n 5 python solution.py --text test_cases/text_4.txt --vocab test_cases/vocab_4.txt --stopwords test_cases/stopwords_4.txt --pattern 4

Output:
Running Pattern 4 with 5 processes.
Term-Frequency (TF) Result:
algorithms: 33
complexity: 3
data: 7
efficient: 1
memory: 2
patterns: 2
problems: 4
search: 2
solutions: 2
sort: 10
time: 3

Document-Frequency (DF) Result:
algorithms: 33
complexity: 3
data: 7
efficient: 1
memory: 2
patterns: 2
problems: 4
search: 2
solutions: 2
sort: 8
time: 3
\end{lstlisting}

%------------------------------------------------------------------------------
\subsection{Test Case 5 Results}
%------------------------------------------------------------------------------

\subsubsection{Pattern 1 (4 processes)}
\begin{lstlisting}
Command: mpiexec -n 4 python solution.py --text test_cases/text_5.txt --vocab test_cases/vocab_5.txt --stopwords test_cases/stopwords_5.txt --pattern 1

Output:
Running Pattern 1 with 4 processes.
Term-Frequency (TF) Result:
ancient: 18
architecture: 3
bosphorus: 3
city: 17
coast: 6
cuisine: 2
culture: 1
istanbul: 2
mosque: 4
mountains: 3
museum: 3
tea: 1
tourism: 1
traditions: 1
turkish: 45
\end{lstlisting}

\subsubsection{Pattern 2 (5 processes)}
\begin{lstlisting}
Command: mpiexec -n 5 python solution.py --text test_cases/text_5.txt --vocab test_cases/vocab_5.txt --stopwords test_cases/stopwords_5.txt --pattern 2

Output:
Running Pattern 2 with 5 processes.
Term-Frequency (TF) Result:
ancient: 18
architecture: 3
bosphorus: 3
city: 17
coast: 6
cuisine: 2
culture: 1
istanbul: 2
mosque: 4
mountains: 3
museum: 3
tea: 1
tourism: 1
traditions: 1
turkish: 45
\end{lstlisting}

\subsubsection{Pattern 3 (9 processes)}
\begin{lstlisting}
Command: mpiexec -n 9 python solution.py --text test_cases/text_5.txt --vocab test_cases/vocab_5.txt --stopwords test_cases/stopwords_5.txt --pattern 3

Output:
Running Pattern 3 with 9 processes.
Term-Frequency (TF) Result:
ancient: 18
architecture: 3
bosphorus: 3
city: 17
coast: 6
cuisine: 2
culture: 1
istanbul: 2
mosque: 4
mountains: 3
museum: 3
tea: 1
tourism: 1
traditions: 1
turkish: 45
\end{lstlisting}

\subsubsection{Pattern 4 (5 processes)}
\begin{lstlisting}
Command: mpiexec -n 5 python solution.py --text test_cases/text_5.txt --vocab test_cases/vocab_5.txt --stopwords test_cases/stopwords_5.txt --pattern 4

Output:
Running Pattern 4 with 5 processes.
Term-Frequency (TF) Result:
ancient: 18
architecture: 3
bosphorus: 3
city: 17
coast: 6
cuisine: 2
culture: 1
istanbul: 2
mosque: 4
mountains: 3
museum: 3
tea: 1
tourism: 1
traditions: 1
turkish: 45

Document-Frequency (DF) Result:
ancient: 18
architecture: 3
bosphorus: 3
city: 17
coast: 6
cuisine: 2
culture: 1
istanbul: 2
mosque: 3
mountains: 3
museum: 3
tea: 1
tourism: 1
traditions: 1
turkish: 45
\end{lstlisting}

%==============================================================================
\section{Work Sharing}
%==============================================================================

% TODO: MODIFY THIS SECTION WITH YOUR ACTUAL CONTRIBUTIONS

\begin{table}[h]
\centering
\begin{tabular}{|l|l|p{8cm}|}
\hline
\textbf{Component} & \textbf{Contributor} & \textbf{Description} \\
\hline
NLP Helper Functions & Student 1 & Implemented \texttt{load\_file}, \texttt{to\_lower}, \texttt{strip\_punct}, \texttt{filter\_stopwords}, \texttt{clean\_text} \\
\hline
TF/DF Computation & Student 1 & Implemented \texttt{compute\_tf} and \texttt{compute\_df} functions. \\
\hline
Pattern 1 & Student 2 & Implemented end-to-end processing. \\
\hline
Pattern 2 & Student 2 & Implemented linear pipeline. \\
\hline
Pattern 3 & Student 1 & Implemented parallel pipelines. \\
\hline
Pattern 4 & Both & Jointly implemented task parallelism. \\
\hline
Test Cases & Student 2 & Created 5 test cases. \\
\hline
Report & Student 1 & Wrote project report. \\
\hline
Testing \& Debugging & Both & Ran all test cases and fixed issues. \\
\hline
\end{tabular}
\caption{Work Distribution}
\end{table}

\subsection{Detailed Contributions}

\textbf{Student 1 Name (Student ID: XXXXXXXXX)}:
\begin{itemize}
    \item Designed and implemented the core NLP preprocessing functions
    \item Implemented Pattern 3 (Parallel Pipelines) including the two-level chunking strategy
    \item Contributed to Pattern 4 implementation, specifically the deadlock prevention mechanism
    \item Wrote the project report and code documentation
    \item Conducted code review and optimization
\end{itemize}

\textbf{Student 2 Name (Student ID: XXXXXXXXX)}:
\begin{itemize}
    \item Implemented Pattern 1 (End-to-End Processing) with balanced workload distribution
    \item Implemented Pattern 2 (Linear Pipeline) with proper termination signal handling
    \item Contributed to Pattern 4 implementation, specifically the TF/DF task splitting
    \item Created all 5 test cases with varying complexity levels
    \item Performed testing and validation of all patterns
\end{itemize}

%==============================================================================
\section{Conclusion}
%==============================================================================

In this project, we implemented four MPI patterns for parallel text processing. We verified that all patterns produce the same correct results. We stuck to the requirement of using only \texttt{MPI\_Send} and \texttt{MPI\_Recv}.

We successfully handled challenges like distributing work evenly and preventing deadlocks. The tests show that our code works correctly for different input sizes and numbers of processes.

\end{document}
